\documentclass[10pt,a4paper]{amsart}
\usepackage{amsmath,amssymb,amsthm,graphicx}
% Retained minimal preamble for the Stacks Project chapter sites-cohomology.tex (Cohomology on Sites).
% Verbatim excerpt of the chapter preamble: Xy-pic setup (native preamble lines 14-18) and the native
% theorem declarations with the mathematical macros (native preamble lines 154-217). Copied mechanically;
% nothing added, removed or redefined. The wrapper class and the amsmath/amssymb/amsthm/graphicx base
% packages are supplied by the assembler, so the class/externaldocument part of the native preamble is omitted.
% For commutative diagrams we use Xy-pic
\usepackage[all]{xy}

% We use 2cell for 2-commutative diagrams.
\xyoption{2cell}
% Theorem environments.
%
\theoremstyle{plain}
\newtheorem{theorem}[subsection]{Theorem}
\newtheorem{proposition}[subsection]{Proposition}
\newtheorem{lemma}[subsection]{Lemma}

\theoremstyle{definition}
\newtheorem{definition}[subsection]{Definition}
\newtheorem{example}[subsection]{Example}
\newtheorem{exercise}[subsection]{Exercise}
\newtheorem{situation}[subsection]{Situation}

\theoremstyle{remark}
\newtheorem{remark}[subsection]{Remark}
\newtheorem{remarks}[subsection]{Remarks}

\numberwithin{equation}{subsection}

% Macros
%
\def\lim{\mathop{\mathrm{lim}}\nolimits}
\def\colim{\mathop{\mathrm{colim}}\nolimits}
\def\Spec{\mathop{\mathrm{Spec}}}
\def\Hom{\mathop{\mathrm{Hom}}\nolimits}
\def\Ext{\mathop{\mathrm{Ext}}\nolimits}
\def\SheafHom{\mathop{\mathcal{H}\!\mathit{om}}\nolimits}
\def\SheafExt{\mathop{\mathcal{E}\!\mathit{xt}}\nolimits}
\def\Sch{\mathit{Sch}}
\def\Mor{\mathop{\mathrm{Mor}}\nolimits}
\def\Ob{\mathop{\mathrm{Ob}}\nolimits}
\def\Sh{\mathop{\mathit{Sh}}\nolimits}
\def\NL{\mathop{N\!L}\nolimits}
\def\CH{\mathop{\mathrm{CH}}\nolimits}
\def\proetale{{pro\text{-}\acute{e}tale}}
\def\etale{{\acute{e}tale}}
\def\QCoh{\mathit{QCoh}}
\def\Ker{\mathop{\mathrm{Ker}}}
\def\Im{\mathop{\mathrm{Im}}}
\def\Coker{\mathop{\mathrm{Coker}}}
\def\Coim{\mathop{\mathrm{Coim}}}

% Boxtimes
%
\DeclareMathSymbol{\boxtimes}{\mathbin}{AMSa}{"02}

%
% Macros for moduli stacks/spaces
%
\def\QCohstack{\mathcal{QC}\!\mathit{oh}}
\def\Cohstack{\mathcal{C}\!\mathit{oh}}
\def\Spacesstack{\mathcal{S}\!\mathit{paces}}
\def\Quotfunctor{\mathrm{Quot}}
\def\Hilbfunctor{\mathrm{Hilb}}
\def\Curvesstack{\mathcal{C}\!\mathit{urves}}
\def\Polarizedstack{\mathcal{P}\!\mathit{olarized}}
\def\Complexesstack{\mathcal{C}\!\mathit{omplexes}}
% \Pic is the operator that assigns to X its picard group, usage \Pic(X)
% \Picardstack_{X/B} denotes the Picard stack of X over B
% \Picardfunctor_{X/B} denotes the Picard functor of X over B
\def\Pic{\mathop{\mathrm{Pic}}\nolimits}
\def\Picardstack{\mathcal{P}\!\mathit{ic}}
\def\Picardfunctor{\mathrm{Pic}}
\def\Deformationcategory{\mathcal{D}\!\mathit{ef}}
\usepackage{cleveref}
\crefname{originalreference}{original reference}{original references}
\Crefname{originalreference}{Original reference}{Original references}
\makeatletter
\newlabel{lemma-compact-in-terms-of-generators}{{094B}{source}{}{sourcepointer-1}{}}
\newlabel{lemma-compact-in-terms-of-generators@cref}{{[originalreference][0][]094B}{[1][0][]source}{}{}{}}
\makeatother
\crefname{originallemma}{original lemma}{original lemmas}
\Crefname{originallemma}{Original lemma}{Original lemmas}
\makeatletter
\newlabel{sites-lemma-directed-colimits-sections}{{0738}{source}{}{sourcepointer-2}{}}
\newlabel{sites-lemma-directed-colimits-sections@cref}{{[originallemma][0][]0738}{[1][0][]source}{}{}{}}
\makeatother
\crefname{originallemma}{original lemma}{original lemmas}
\Crefname{originallemma}{Original lemma}{Original lemmas}
\makeatletter
\newlabel{sites-modules-lemma-module-quotient-flat}{{03EW}{source}{}{sourcepointer-3}{}}
\newlabel{sites-modules-lemma-module-quotient-flat@cref}{{[originallemma][0][]03EW}{[1][0][]source}{}{}{}}
\makeatother
\title{If every object has a cover by quasi-compact objects then compact objects are summands of finite complexes of extension sheaves}
\author{The Stacks Project Authors}
\date{stacks-project:sites-cohomology.tex:local-snapshot-20261011; GNU FDL 1.2 or later}
\begin{document}
\maketitle
\noindent\textit{Editorial note.} All delivered mathematical text is the author's own, copied from the pinned Stacks Project chapter file sites-cohomology.tex; no statement or proof step was rewritten, reordered or removed.\par\medskip
\noindent\textit{Editorial note.} The counted proof is the author's complete argument: the proof is written out in full: the choice of the generating set, and the verification of the two conditions of the criterion with the reason for each.\par\medskip
\noindent\textit{Editorial note.} The label lemma-compact-objects-if-enough-qc names the counted statement, and the proof environment immediately following it in the source is delivered verbatim as its proof.\par\medskip
\noindent\textit{Editorial note.} The counted statement is carried by the author's own lemma environment in the chapter source.\par\medskip
\noindent\textit{Editorial note.} The statement and proof are copied from the pinned chapter file; the three cited inputs are kept as external pointers with their canonical Stacks tags.\par\medskip
\noindent\textit{Allowed standard prerequisites (editorial).}\begin{itemize}
\item Compactness in terms of a generating set. Same chapter, the lemma describing compact objects through a set of generators satisfying two conditions. Supplies the criterion applied to the set of extension sheaves j\_!O\_U over quasi-compact U.
\item Every module is a quotient of a direct sum of extension sheaves. Modules on Sites, the lemma representing an O-module as a quotient of a direct sum of j\_!O\_U. Supplies condition (1) of the criterion, together with the covering assumption.
\item Sections commute with directed colimits. Sites, the lemma computing sections of a directed colimit as the colimit of the sections. Supplies condition (2) of the criterion through the identity Hom(j\_!O\_U, F) = F(U).
\end{itemize}\medskip
\ifdefined\hypertarget\hypertarget{sourcepointer-1}{}\fi
\noindent\textit{Source reference.} Original reference 094B: sources/stacks-sites-cohomology/pinned-original-source.tex, source lines 13934--13934. Prerequisite pointer invoked at proof line(s) 14045; it names a standard fact of the same chapter that the counted argument uses, and every step of the counted proof built on it is retained verbatim in the delivered body.\par\medskip
\ifdefined\hypertarget\hypertarget{sourcepointer-2}{}\fi
\noindent\textit{Source reference.} Original lemma 0738: sources/stacks-chow/tags.txt, source lines 9052--9052. Prerequisite pointer invoked at proof line(s) 14056; it names a standard fact of a sibling chapter that the counted argument uses, and every step of the counted proof built on it is retained verbatim in the delivered body.\par\medskip
\ifdefined\hypertarget\hypertarget{sourcepointer-3}{}\fi
\noindent\textit{Source reference.} Original lemma 03EW: sources/stacks-chow/tags.txt, source lines 4471--4471. Prerequisite pointer invoked at proof line(s) 14049; it names a standard fact of a sibling chapter that the counted argument uses, and every step of the counted proof built on it is retained verbatim in the delivered body.\par\medskip
\begin{lemma}
\label{lemma-compact-objects-if-enough-qc}
Let $(\mathcal{C}, \mathcal{O})$ be a ringed site. Assume every object of
$\mathcal{C}$ has a covering by quasi-compact objects. Then every
compact object of $D(\mathcal{O})$ is a direct summand in $D(\mathcal{O})$
of a finite complex whose terms are finite direct sums of
$\mathcal{O}$-modules of the form $j_!\mathcal{O}_U$
where $U$ is a quasi-compact object of $\mathcal{C}$.
\end{lemma}

\begin{proof}
Apply Lemma \ref{lemma-compact-in-terms-of-generators}
where $S \subset \Ob(\textit{Mod}(\mathcal{O}))$ is the set of modules
of the form $j_!\mathcal{O}_U$ with $U \in \Ob(\mathcal{C})$
quasi-compact. Assumption (1) holds by
Modules on Sites, Lemma \ref{sites-modules-lemma-module-quotient-flat}
and the assumption that every $U$ can be covered by quasi-compact
objects. Assumption (2) follows as
$$
\Hom_\mathcal{O}(j_!\mathcal{O}_U, \mathcal{F}) = \mathcal{F}(U)
$$
which commutes with direct sums by
Sites, Lemma \ref{sites-lemma-directed-colimits-sections}.
\end{proof}

\end{document}
